%==============================================================================
%  Rapport : Energie et planting d'arbres
%  Compilation : pdflatex (deux fois, pour les references et la table des matieres)
%
%  IMPORTANT - Emplacement des images
%  ----------------------------------
%  Les 16 images se trouvent dans le dossier Windows :
%      C:\Users\LENOVO\OneDrive\PLANTING
%  En LaTeX, on ecrit le chemin avec des barres obliques " / " (et non " \ ").
%  Le chemin est declare une seule fois ci-dessous avec \graphicspath : il suffit
%  ensuite d'appeler chaque image par son nom de fichier.
%
%  Les noms d'images utilises ici sont photo1.jpg, photo2.jpg, ... photo16.jpg.
%  -> Renomme tes 16 fichiers ainsi, OU remplace les noms dans le document.
%  Formats acceptes par pdflatex : .jpg, .png, .pdf
%==============================================================================

\documentclass[12pt,a4paper]{article}

%------------------------------------------------------------------ Encodage et langue
\usepackage[utf8]{inputenc}      % saisie des caracteres accentues
\usepackage[T1]{fontenc}         % bonne gestion des accents a l'affichage
\usepackage{lmodern}             % police vectorielle nette
\usepackage[french]{babel}       % regles typographiques francaises

%------------------------------------------------------------------ Mise en page
\usepackage[a4paper,margin=2.5cm]{geometry}  % marges
\usepackage{setspace}            % interligne
\onehalfspacing                  % interligne 1,5
\usepackage{microtype}           % justification plus propre

%------------------------------------------------------------------ Images et tableaux
\usepackage{graphicx}            % insertion des images
\usepackage{subcaption}          % plusieurs images cote a cote
\usepackage{float}               % option [H] pour fixer la position des figures
\usepackage{booktabs}            % tableaux soignes

%------------------------------------------------------------------ Couleurs et liens
\usepackage{xcolor}
\definecolor{vertfeuille}{RGB}{34,110,45}   % vert pour les titres
\usepackage[colorlinks=true,linkcolor=vertfeuille,urlcolor=vertfeuille]{hyperref}

%------------------------------------------------------------------ Titres de section en vert
\usepackage{titlesec}
\titleformat{\section}{\Large\bfseries\color{vertfeuille}}{\thesection}{1em}{}
\titleformat{\subsection}{\large\bfseries\color{vertfeuille}}{\thesubsection}{1em}{}

%------------------------------------------------------------------ En-tete et pied de page
\usepackage{fancyhdr}
\pagestyle{fancy}
\fancyhf{}
\fancyhead[L]{\small Energie et planting d'arbres}
\fancyhead[R]{\small\leftmark}
\fancyfoot[C]{\small\thepage}
\renewcommand{\headrulewidth}{0.4pt}

%------------------------------------------------------------------ CHEMIN DES IMAGES
% Toutes les images sont cherchees dans ce dossier (barres obliques " / ").
% La barre finale est obligatoire.
\graphicspath{{C:/Users/LENOVO/OneDrive/PLANTING/}}

%------------------------------------------------------------------ Commande pratique
% \photo{nom_du_fichier}{legende} insere une image centree, large de 80 % du texte.
\newcommand{\photo}[2]{%
  \begin{figure}[H]
    \centering
    \includegraphics[width=0.8\textwidth]{#1}
    \caption{#2}
  \end{figure}}

%==============================================================================
\begin{document}

%------------------------------------------------------------------ PAGE DE TITRE
\begin{titlepage}
  \centering
  \vspace*{2cm}
  {\Huge\bfseries\color{vertfeuille} Energie et planting d'arbres\par}
  \vspace{0.5cm}
  {\Large Vers un environnement durable\par}
  \vspace{2cm}

  % Image de couverture (la premiere photo).
  \includegraphics[width=0.7\textwidth]{photo1.jpg}

  \vfill
  {\large Presente par : \textbf{KOUAME Koffi Fidele}\par}
  \vspace{0.3cm}
  {\large Date : \today\par}
\end{titlepage}

%------------------------------------------------------------------ TABLE DES MATIERES
\tableofcontents
\newpage

%==============================================================================
\section{Introduction}
%==============================================================================
Le rechauffement climatique et la demande croissante en energie poussent a
rechercher des solutions durables. Le planting d'arbres est l'une des reponses
les plus accessibles : en absorbant le dioxyde de carbone, les arbres limitent
l'effet de serre tout en soutenant les besoins energetiques des populations,
notamment a travers le bois-energie gere de facon responsable.

Ce rapport presente une activite de planting d'arbres, ses objectifs, son
deroulement et ses resultats, illustres par une serie de photographies.

%==============================================================================
\section{Contexte et objectifs}
%==============================================================================

\subsection{Lien entre energie et arbres}
Les arbres jouent un double role. D'une part, ils stockent le carbone et
regulent le climat local. D'autre part, ils constituent une source d'energie
renouvelable lorsqu'ils sont exploites de maniere durable (bois de chauffe,
charbon de bois ameliore, biomasse). Planter des arbres, c'est donc investir a
la fois dans le climat et dans l'energie de demain.

\subsection{Objectifs de l'activite}
\begin{itemize}
  \item Sensibiliser la communaute a la protection de l'environnement.
  \item Contribuer a la reduction des emissions de gaz a effet de serre.
  \item Mettre en place une ressource durable de bois-energie.
  \item Restaurer les sols et preserver la biodiversite.
\end{itemize}

%------------------------------------------------------------------ Photo de contexte
\photo{photo2.jpg}{Site choisi pour le planting avant le debut des travaux.}

%==============================================================================
\section{Materiel et methode}
%==============================================================================

\subsection{Materiel utilise}
\begin{table}[H]
  \centering
  \begin{tabular}{@{}ll@{}}
    \toprule
    \textbf{Materiel} & \textbf{Usage} \\
    \midrule
    Plants (jeunes arbres)   & Mise en terre \\
    Beches et pelles         & Creusage des trous \\
    Arrosoirs                & Arrosage apres plantation \\
    Gants                    & Protection des mains \\
    Piquets et cordes        & Alignement des plants \\
    \bottomrule
  \end{tabular}
  \caption{Materiel mobilise pour l'activite.}
\end{table}

\subsection{Etapes du planting}
\begin{enumerate}
  \item Preparation et nettoyage du terrain.
  \item Marquage des emplacements a intervalles reguliers.
  \item Creusage des trous de plantation.
  \item Mise en terre des plants et rebouchage.
  \item Arrosage et protection des jeunes arbres.
\end{enumerate}

%------------------------------------------------------------------ Photos de preparation
\photo{photo3.jpg}{Preparation du terrain par l'equipe.}
\photo{photo4.jpg}{Marquage et alignement des emplacements.}

%==============================================================================
\section{Deroulement de l'activite}
%==============================================================================
Cette section rassemble les photographies prises pendant le planting. Les images
sont presentees deux par deux pour faciliter la lecture.

%------------------------------------------------------------------ GALERIE 2 PAR 2
% Chaque bloc "figure" contient deux images cote a cote (deux "subfigure").
% Pour ajouter ou retirer une image, copie ou supprime un bloc "subfigure".

\begin{figure}[H]
  \centering
  \begin{subfigure}{0.48\textwidth}
    \centering\includegraphics[width=\textwidth]{photo5.jpg}
    \caption{Creusage des trous.}
  \end{subfigure}\hfill
  \begin{subfigure}{0.48\textwidth}
    \centering\includegraphics[width=\textwidth]{photo6.jpg}
    \caption{Mise en terre d'un plant.}
  \end{subfigure}
  \caption{Les premieres etapes de la plantation.}
\end{figure}

\begin{figure}[H]
  \centering
  \begin{subfigure}{0.48\textwidth}
    \centering\includegraphics[width=\textwidth]{photo7.jpg}
    \caption{Rebouchage et tassement du sol.}
  \end{subfigure}\hfill
  \begin{subfigure}{0.48\textwidth}
    \centering\includegraphics[width=\textwidth]{photo8.jpg}
    \caption{Arrosage des jeunes arbres.}
  \end{subfigure}
  \caption{La mise en terre et l'arrosage.}
\end{figure}

\begin{figure}[H]
  \centering
  \begin{subfigure}{0.48\textwidth}
    \centering\includegraphics[width=\textwidth]{photo9.jpg}
    \caption{Travail en equipe.}
  \end{subfigure}\hfill
  \begin{subfigure}{0.48\textwidth}
    \centering\includegraphics[width=\textwidth]{photo10.jpg}
    \caption{Participation de la communaute.}
  \end{subfigure}
  \caption{La mobilisation des participants.}
\end{figure}

\begin{figure}[H]
  \centering
  \begin{subfigure}{0.48\textwidth}
    \centering\includegraphics[width=\textwidth]{photo11.jpg}
    \caption{Plant protege par un tuteur.}
  \end{subfigure}\hfill
  \begin{subfigure}{0.48\textwidth}
    \centering\includegraphics[width=\textwidth]{photo12.jpg}
    \caption{Rangee d'arbres fraichement plantes.}
  \end{subfigure}
  \caption{Les arbres mis en place.}
\end{figure}

\begin{figure}[H]
  \centering
  \begin{subfigure}{0.48\textwidth}
    \centering\includegraphics[width=\textwidth]{photo13.jpg}
    \caption{Vue d'ensemble du site.}
  \end{subfigure}\hfill
  \begin{subfigure}{0.48\textwidth}
    \centering\includegraphics[width=\textwidth]{photo14.jpg}
    \caption{Suivi et entretien.}
  \end{subfigure}
  \caption{Le site apres la plantation.}
\end{figure}

\begin{figure}[H]
  \centering
  \begin{subfigure}{0.48\textwidth}
    \centering\includegraphics[width=\textwidth]{photo15.jpg}
    \caption{Photo de groupe des participants.}
  \end{subfigure}\hfill
  \begin{subfigure}{0.48\textwidth}
    \centering\includegraphics[width=\textwidth]{photo16.jpg}
    \caption{Les arbres, source d'energie durable.}
  \end{subfigure}
  \caption{La cloture de l'activite.}
\end{figure}

%==============================================================================
\section{Resultats et impact}
%==============================================================================
\begin{itemize}
  \item \textbf{Nombre d'arbres plantes :} 16 (a adapter selon ton bilan reel).
  \item \textbf{Superficie couverte :} a completer.
  \item \textbf{Participants mobilises :} a completer.
\end{itemize}

A terme, les arbres plantes contribueront a capter le dioxyde de carbone,
a ameliorer la qualite de l'air et a fournir une ressource de bois-energie
renouvelable, a condition d'un entretien regulier.

%==============================================================================
\section{Conclusion}
%==============================================================================
Le planting d'arbres relie directement la protection de l'environnement et la
question de l'energie. Chaque arbre plante est un pas vers un cadre de vie plus
sain et une energie plus durable. La reussite du projet repose desormais sur le
suivi et l'entretien des plants dans la duree.

\end{document}
